
    \centering
    % ESTA LÍNEA es la que crea el PDF en la carpeta /tikz
    \tikzsetnextfilename{ejercicio21_figura} 
    % Colores originales
    \definecolor{zzttqq}{rgb}{0.6,0.2,0}
    \definecolor{ududff}{rgb}{0.30196,0.30196,1}

    \begin{tikzpicture}[line cap=round, line join=round]
    \begin{axis}[
    view={110}{25}, % Ángulo de cámara para perspectiva 3D (azimut y elevación)
    axis lines=center,
    xlabel={$x$}, ylabel={$y$}, zlabel={$z$},
    xmin=-2, xmax=10,
    ymin=-2, ymax=10,
    zmin=-2, zmax=10,
    xtick={-2,2,4,6,8},
    ytick={-2,2,4,6,8},
    ztick={-2,2,4,6,8},
    ticklabel style={font=\scriptsize},
    width=10cm, height=10cm
    ]

    % 1. Definimos los vértices del paralelepípedo en 3D real
    \coordinate (P000) at (0,0,0);
    \coordinate (P100) at (5,2,1); % Vector u
    \coordinate (P010) at (2,4,2); % Vector v
    \coordinate (P001) at (1,3,6); % Vector w

    \coordinate (P110) at (7,6,3); % u + v
    \coordinate (P101) at (6,5,7); % u + w
    \coordinate (P011) at (3,7,8); % v + w
    \coordinate (P111) at (8,9,9); % u + v + w

    % 2. Dibujamos las aristas traseras (ocultas) punteadas
    \draw[very thick, dashed, color=zzttqq] (P000) -- (P010);
    \draw[very thick, dashed, color=zzttqq] (P010) -- (P110);
    \draw[very thick, dashed, color=zzttqq] (P010) -- (P011);

    % 3. Dibujamos las caras visibles (con relleno y opacidad)
    % Cara frontal-derecha
    \filldraw[fill=zzttqq, fill opacity=0.15, draw=zzttqq, very thick] 
    (P000) -- (P100) -- (P101) -- (P001) -- cycle;
    % Cara lateral-derecha
    \filldraw[fill=zzttqq, fill opacity=0.15, draw=zzttqq, very thick] 
    (P100) -- (P110) -- (P111) -- (P101) -- cycle;
    % Cara superior
    \filldraw[fill=zzttqq, fill opacity=0.15, draw=zzttqq, very thick] 
    (P001) -- (P101) -- (P111) -- (P011) -- cycle;

    % 4. Añadimos flechas a los vectores principales
    \draw[very thick, ->, >=Triangle, color=zzttqq] (P000) -- (P100);
    \draw[very thick, dashed, ->, >=Triangle, color=zzttqq] (P000) -- (P010);
    \draw[very thick, ->, >=Triangle, color=zzttqq] (P000) -- (P001);

    % 5. Dibujamos los puntos y etiquetas
    \fill[ududff] (P000) circle (2.5pt); 
    \fill[black!60] (P100) circle (2.5pt); 
    \fill[black!60] (P010) circle (2.5pt); 
    \fill[black!60] (P001) circle (2.5pt); 
    \fill[black!60] (P110) circle (2.5pt); 
    \fill[black!60] (P101) circle (2.5pt); 
    \fill[black!60] (P011) circle (2.5pt); 
    \fill[black!60] (P111) circle (2.5pt); 

    \end{axis}
    \end{tikzpicture}
    \caption{Representación de la región pedida.  \href{https://www.geogebra.org/m/bjnjeat8}{Ver en GeoGebra}}
